% ============================================================
%  Função Afim — Apostila Premium | PratiqueJá
%  Engine: XeLaTeX
% ============================================================
\documentclass[12pt]{article}
\usepackage{../../pratiqueja}
\usepackage{tikz}

% \hlm — destaque de resultado em modo-math (cor sem bold)
\newcommand{\hlm}[1]{{\color{darkblue}#1}}

\setsubject{Função Afim}

\begin{document}
\pagenumbering{arabic}
\setstretch{1.2}
\color{bodytext}

\heroheader{Função Afim}%
           {Coeficientes, Zero, Estudo do Sinal e Aplicações}%
           {Matemática · Avançado}

\setlength{\columnsep}{18pt}
\setlength{\columnseprule}{0.5pt}
\def\columnseprulecolor{\color{exborder}}

\begin{multicols}{2}

%─────────────────────────────────────────────────────────────
\TW{Def}{badgepurple}{Definição}{O que é a função afim?}

\regra{Função $f:\mathbb{R}\to\mathbb{R}$ dada por $f(x) = ax + b$, com
$a,b \in \mathbb{R}$ e $\mathbf{a \neq 0}$. No plano é uma \textbf{reta
não horizontal} ($y = ax + b$) — também dita \textbf{função de 1º grau}.}

\obs{$a$ = coeficiente \textbf{angular} (inclinação / taxa de variação);
$b$ = coeficiente \textbf{linear} $= f(0)$ (onde corta o eixo $y$).}

\exemp{%
  $f(x) = 3x + 4$: \ $a = 3,\ b = 4$\\
  $f(2) = 3\cdot2 + 4 = \hlm{10}$%
}

%─────────────────────────────────────────────────────────────
\TW{a}{darkblue}{Coeficiente Angular}{Taxa de variação e inclinação}

\regra{A partir de dois pontos $A,B$ da reta:}
\formula{$a = \dfrac{\Delta y}{\Delta x} = \dfrac{y_B - y_A}{x_B - x_A}$}

\exemp{%
  $A=(2,3),\ B=(3,5)$: $a = \dfrac{5-3}{3-2} = \hlm{2}$\\
  {\footnotesize $+1$ em $x$ → $+2$ em $y$ (crescente)}\\[3pt]
  $A=(1,4),\ B=(3,0)$: $a = \dfrac{0-4}{3-1} = \hlm{-2}$\\
  {\footnotesize $+1$ em $x$ → $-2$ em $y$ (decrescente)}%
}
\obs{$a > 0$ → função \textbf{crescente}; \ $a < 0$ → \textbf{decrescente}.}

%─────────────────────────────────────────────────────────────
\TW{b}{badgegreen}{Coeficiente Linear}{Interseção com o eixo $y$}

\regra{$b = f(0)$ — ponto onde a reta corta o eixo $y$. Conhecidos $a$ e
um ponto $(x,y)$:}
\formula{$b = y - a\cdot x$}

\exemp{%
  Reta por $(2,7)$ com $a=3$:\\
  $b = 7 - 3\cdot2 = \hlm{1} \Rightarrow f(x) = 3x + 1$%
}

\SH{Determinar a função pelo gráfico}
\begin{center}\vspace{2pt}
\begin{tikzpicture}[scale=0.33,font=\scriptsize]
\draw[help lines, gray!30] (-6,-5) grid (6,7);
\draw[->,line width=0.8pt] (-6,0)--(6.6,0) node[right]{$x$};
\draw[->,line width=0.8pt] (0,-5)--(0,7.6) node[above]{$y$};
\draw[darkblue,line width=1.2pt] (-6,-5)--(6,7);
\fill[vered] (-5,-4) circle (5pt);
\node[vered,anchor=west] at (-4.4,-4){$(-5,-4)$};
\fill[vered] (5,6) circle (5pt);
\node[vered,anchor=east] at (4.4,6){$(5,6)$};
\end{tikzpicture}
\end{center}

\exemp{%
  Por $(-5,-4)$ e $(5,6)$:\\
  $a = \dfrac{6-(-4)}{5-(-5)} = \dfrac{10}{10} = 1$\\
  $6 = 1\cdot5 + b \Rightarrow b = 1$\\
  $\Rightarrow \hlm{f(x) = x + 1}$%
}

%─────────────────────────────────────────────────────────────
\TW{$x_0$}{badgepurple}{Zero da Função}{Onde a reta cruza o eixo $x$}

\regra{Valor de $x$ tal que $f(x) = 0$:}
\formula{$ax + b = 0 \;\Rightarrow\; x_0 = -\dfrac{b}{a}$}
\obs{Existe e é \textbf{único} para toda afim ($a \neq 0$): a reta cruza
o eixo $x$ em exatamente um ponto.}

\exemp{%
  $2x - 6 = 0 \Rightarrow x_0 = \hlm{3}$\\[2pt]
  $-3x + 9 = 0 \Rightarrow x_0 = \hlm{3}$\\[2pt]
  $x + 1 = 0 \Rightarrow x_0 = \hlm{-1}$%
}

%─────────────────────────────────────────────────────────────
\TW{$\pm$}{darkblue}{Estudo do Sinal}{Quando $f$ é positiva, nula ou negativa}

\regra{O sinal de $f$ \textbf{muda no zero} $x_0$ e depende do sinal de
$a$.}

\exemp{%
  \textbf{$a>0$ (crescente):}\\
  $f(x) < 0$ se $x < x_0$; \quad $f(x) > 0$ se $x > x_0$\\[3pt]
  \textbf{$a<0$ (decrescente):}\\
  $f(x) > 0$ se $x < x_0$; \quad $f(x) < 0$ se $x > x_0$%
}

\SH{Inequações}
\exemp{%
  $2x - 6 > 0$: \ $x_0=3,\ a>0$ → \hlm{$x > 3$}\\[2pt]
  $-3x + 9 \geq 0$: \ $x_0=3,\ a<0$ → \hlm{$x \leq 3$}%
}

%─────────────────────────────────────────────────────────────
\TW{Esp}{badgegreen}{Casos Particulares}{Função linear e função constante}

\exemp{%
  \textbf{Linear} $(b=0)$: $f(x) = ax$ passa pela \textbf{origem}.\\
  Ex.: $f(x) = 5x$ → zero em $x_0 = 0$%
}
\exemp{%
  \textbf{Constante} $(a=0)$: $f(x) = b$, reta \textbf{horizontal}.\\
  Ex.: $f(x) = 4$ → não cruza o eixo $x$ (sem zero)%
}
\nota{Toda função linear é afim, mas nem toda afim é linear. A
\textbf{constante não é afim} (pois $a=0$).}

%─────────────────────────────────────────────────────────────
\TW{Prob}{badgepurple}{Problemas Contextualizados}{Aplicação em situações reais}

\regra{\textbf{Estratégia:} (1) achar a taxa $a$ e o valor inicial $b$;
(2) montar $f(x) = ax + b$; (3) responder.}

\SH{P1 — Táxi}
\exemp{%
  Bandeirada R\$\,$5$ $+$ R\$\,$2{,}50$/km: \ $f(x) = 2{,}50x + 5$\\
  $f(x) > 30$: \ $2{,}50x > 25 \Rightarrow \hlm{x > 10\text{ km}}$%
}

\SH{P2 — Temperatura}
\exemp{%
  Parte de $80°$C, cai $5°$C/min: \ $T(x) = -5x + 80$\\
  $-5x + 80 = 30 \Rightarrow -5x = -50 \Rightarrow \hlm{x = 10\text{ min}}$%
}

\end{multicols}

\end{document}
